\section{Construction}\label{sec:agg-creds-construction}

The construction assumes that the registration authority checks that users'
public keys are unique, and that users hold the corresponding secret keys.
This prevents users from registering others' public keys and pooling credentials
across identities.  The registration authority assigns each user exactly one
tag, and signs only messages of the prescribed form: public keys and tags. We assume that verifiers
accept credentials only from the set of issuers with published public keys. This
prevents a malicious user from being able to impersonate an issuer to an honest
verifier.

Each party's responsibilities fix the building block it needs. The
registration authority only links a public key to a tag, so an existentially
unforgeable signature scheme is enough. Presentations should show one short
object, so issuers' signatures must aggregate. Ordinary aggregate signatures
would let two users pool their credentials, and multi-signatures require every
signer to sign the same message, while issuers sign different attributes.
Tag-based aggregation combines only signatures under one tag, without issuers
coordinating. The proof system must be zero knowledge, so that presentations
reveal only the disclosed values, and simulation extractable, since the security
proof simulates honest users' presentations while extracting from the
adversary's. Each presentation uses the verifier's message as its label, so
it cannot be replayed elsewhere.

\Cref{fig:agg-creds-seq} shows the messages and local steps of the
construction.

\begin{figure}[htbp] \centering \small
\begin{tikzpicture}[seq]
  \foreach \x/\n/\e in {0/RA/4.4, 1/User/13.4, 2/Issuer $i$/10.4, 3/Verifier/14.4} {
    \node[seqhead] at (\x, 0) {\n};
    \draw[seqlife] (\x, 0.6) -- (\x, \e);
  }
  \draw[seqmsg] (0, 1) -- node[seqlabel] {$\rand$} (1, 1);
  \draw[seqmsg] (1, 2) -- node[seqlabel] {$\pk$, proof on $\rand$} (0, 2);
  \node[seqloc] at (0, 3) {new $\Tag$, $\DSsign$};
  \draw[seqmsg] (0, 4) -- node[seqlabel] {$\Tag$, $\psi$} (1, 4);
  \draw[seqmsg] (1, 5) -- node[seqlabel] {$\Tag$, $\pk$, $\psi$, $\atv{i}$} (2, 5);
  \node[seqloc] at (2, 6) {$\DSverify$};
  \draw[seqmsg] (2, 7) -- node[seqlabel] {$\rand$} (1, 7);
  \draw[seqmsg] (1, 8) -- node[seqlabel] {proof on $\rand$} (2, 8);
  \node[seqloc] at (2, 9) {$\ATSsign$};
  \draw[seqmsg] (2, 10) -- node[seqlabel] {$\signature_i$} (1, 10);
  \draw[seqmsg] (3, 11) -- node[seqlabel] {$\msg$} (1, 11);
  \node[seqloc] at (1, 12) {$\AACaggregate$, $\AACpresent$};
  \draw[seqmsg] (1, 13) -- node[seqlabel] {$(\atv{i})_{i \in \ISSUERS}$, $\algoproof$} (3, 13);
  \node[seqloc] at (3, 14) {$\AACverify$};
\end{tikzpicture}
\caption[Aggregate credentials: messages]{Messages and local steps of
registration, issuance from issuer $i$, aggregation and presentation, between
the registration authority, the user, issuer $i$ and a verifier. The proofs on
$\rand$ prove knowledge of $\sk$.} \label{fig:agg-creds-seq} \end{figure}

\paragraph{Setup.} During setup, the registration authority generates a key pair
for the digital signature scheme and public parameters for the proofs used
throughout the protocol. Each issuer independently generates a key pair for the
aggregate tag-based signature scheme. The registration authority's public key,
each issuer's public key, and the public parameters are published. Each issuer
publishes its public key together with a proof of knowledge of its secret key,
and verifiers only accept issuer keys whose proof verifies. The public
parameters $\crs$ consist of the registration authority's public key, the issuer
public keys $\ipk{1}, \ldots, \ipk{n}$ and the parameters of the proof system.

\paragraph{Registration.} Users' keys are related by a one-way function $f$. The user samples $\sk$
uniformly, sets $\pk = f(\sk)$, and sends $\pk$ to the registration authority together with a proof of knowledge of $\sk$ on a
fresh message from the registration authority.
If the proof verifies, the registration authority assigns the user a unique tag
$\Tag$ and signs $(\pk, \Tag)$ with $\DS$, giving $\psi$.

The user initialises their credential state as $(\Tag, \sk, \pk, \psi, \cred_1,
\ldots, \cred_n)$, with $\cred_i = \varnothing$ for all $i \in [1,n]$, and
$\regcred = (\Tag, \pk, \psi)$ is their registration credential.

\paragraph{Issuance.} On input $(\Tag, \pk, \psi, \atv{i})$ from the user, the issuer checks that
$\psi$ verifies, and rejects if it does not.  Otherwise, the user next proves
knowledge of $\sk$.  If this succeeds and the user holds attribute value
$\atv{i}$, the issuer sends back $\signature_i \sample \ATSsign(\isk{i},
\Tag, \atv{i})$. The user sets $\cred_i = (\atv{i}, \signature_i)$.

\paragraph{Aggregation.} Before presenting credentials, the user aggregates the
signatures corresponding to the attributes they wish to disclose into a single
aggregate signature.

\paragraph{Presentation.} The user has $(\Tag, \sk, \pk, \psi, \cred_1, \ldots, \cred_n)$, where $\cred_i
= (\atv{i}, \signature_i)$.  To prove possession of attribute values
$(\atv{i})_{i \in \ISSUERS}$ to a verifier providing a fresh message $\msg$, the
user first constructs the aggregate signature as follows:
%
\begin{equation*}
%
  \signature \leftarrow \ATSaggregate((\signature_i)_{i \in \ISSUERS}),
 %
\end{equation*}
%
which, with the tag and the disclosed values, forms the aggregate credential
$\aggcred$. They then construct a proof of knowledge, with label $\msg$, with the
following parameters:
%
\begin{equation*}
%
  \inst_1 = (\rapk, \ISSUERS, \{\ipk{i}, \atv{i}\}_{i \in \ISSUERS}), \qquad
\witness_1 = (\Tag, \sk, \pk, \psi, \signature).
%
\end{equation*}
%
The relation is the following:
%
\begin{equation}\label{eq:agg-creds-proof-abstract} \begin{gathered}
%
  \rel_1 = \{\ \pk = f(\sk) \ \wedge \
\DSverify(\rapk, (\pk, \Tag), \psi) = 1 \\
%
  \wedge\ \ATSverifyagg((\ipk{i})_{i \in \ISSUERS}, \Tag, (\atv{i})_{i \in
\ISSUERS}, \signature) = 1\ \}.
%
\end{gathered} \end{equation}

The user proves that they possess a valid registration signature, an aggregate
credential under a single tag, and the secret key corresponding to the
registered public key for that tag, without revealing either the tag or the
public key. The aggregate signature is on the disclosed values, one for each
issuer in $\ISSUERS$. The statement contains only $\rapk$, the issuer set, the
issuer public keys and the disclosed values, so the presentation links neither
to the user nor to other presentations.

\paragraph{Verification.} On input $(\{\atv{i}\}_{i \in \ISSUERS}, \msg, \algoproof)$, the verifier
returns the output of the proof verification algorithm.

\begin{theorem}\label{thm:agg-creds-construction-security} If $\DS$ is an
existentially unforgeable digital signature scheme, $\ATS$ an existentially
unforgeable aggregate tag-based signature scheme under tag-based aggregation,
the proof system is zero knowledge and simulation extractable, $f$ is one way,
and the proofs of knowledge of users' and issuers' secret keys are
zero-knowledge proofs of knowledge, then the above protocol securely realises
ideal functionality~$\idfu{\AAC}$.  \end{theorem}

We prove \Cref{thm:agg-creds-construction-security} in \Cref{sec:agg-creds-proof-1}.


\paragraph{Extensions.} The construction extends in several ways.

To bind each presentation to a verifier and a session, $\msg$ contains the
verifier's identity and a fresh nonce. Aggregation is over $\ISSUERS$ only, so
users can present any subset of their credentials by aggregating just that
subset, with no further contact with the issuers or the registration authority.

New issuers can join the system without impacting the construction. The
aggregation of credentials and their presentation is agnostic to the number of
issuers.

An issuer certifying multiple attribute types can hold a different key for each
type. To support an issuer certifying multiple attribute types under one key,
the issuer instead signs $(\Tag, \att{} \| \atv{})$ rather than $(\Tag,
\atv{})$, with $\att{}$ the attribute type and $\atv{}$ the attribute value.

The construction extends to epoch-based revocation. The registration authority
signs triple $(\pk, \Tag, \epoch)$ rather than pair $(\pk, \Tag)$, where
$\epoch$ identifies the current epoch, and issuers sign $(\Tag, \epoch \|
\atv{})$ rather than $(\Tag, \atv{})$, tying each attribute credential to that
epoch. At presentation, users also prove that their registration
signature and aggregate signatures all correspond to the current epoch.

We leave the study of accumulator-based revocation and the revocation of
individual attribute credentials to future work.
